\subsection{Reporte de calidad de los datos}

\subsubsection*{Calidad de los datos de BCS}

\begin{itemize}
  \item De 761 archivos, 136 son copias dentro de la misma vaca, 122 en las carpetas del 14 de octubre y 14 en la carpeta de la vaca 91579 de la primera visita. En las carpetas del 14 de octubre, 13 vacas se descargaron dos veces con exactamente las mismas fotos y hay 18 archivos con nombre tipo ``IMG\_0530(1).HEIC''. En la primera visita hay 53 nombres con ``(1)'' y muchos de ellos son tomas distintas y no copias, por eso las copias se cuentan por huella y no por nombre. Ninguna foto del 29 de septiembre repite una foto del 14 de octubre.
  \item Las vacas 1548 y 2107 tienen 4 fotos iguales y no se sabe a cuál de las dos corresponde cada una. La 2107 no tiene más fotos.
  \item La vaca 2123 aparece dos veces en la bitácora con dos carpetas de fotos distintas, así que no se sabe qué fotos van con cada evaluación.
  \item Tres fotos de la carpeta de la vaca 1538 se llaman ``1538 y 8744\_atras'', por lo que no se sabe a cuál de las dos vacas corresponden.
  \item La carpeta de la vaca 1578 de la primera visita está vacía, aunque la vaca sí tiene fotos del 14 de octubre.
  \item De las 182 fotos JPG, 70 vienen con una orientación EXIF distinta de 1 y se verían giradas si el programa ignora esa etiqueta. Son 20 de las 24 fotos del Samsung S22 y 50 de las 117 fotos JPEG del iPhone 17 Pro Max de la primera visita.
  \item Hay dos formatos de imagen (HEIC y JPG) y 12 resoluciones distintas, y varias fotos de la primera visita son horizontales.
  \item La bitácora tiene 14 horas marcadas con ``x'' y 2 con otro formato, y la hoja no registra la hora. La hora no se usará en el modelo.
  \item Solo 20 de las 79 vacas con BCS están en el conjunto de datos de VMS 3.
  \item ``BCS ayer'' está vacía en las 24 vacas, mientras que ``Evento BCS'' tiene valores parecidos a un BCS. No sabemos si se relacionan ni qué significan.
  \item ``Evento BCS'' no coincide con el BCS de los expertos. De 12 vacas con ambos valores en la bitácora, el conjunto de datos es mayor en 9, igual en 2 y menor en 1. Las mayores diferencias son la vaca 8744 (3.0 contra 4.17), la 8799 (2.5 contra 3.5) y la 1547 (3.25 contra 3.98). No lo usaremos como etiqueta hasta aclarar qué mide y de qué fecha es.
  \item Dos vacas de la primera visita (1549 y 1572) tienen fotos pero no tienen BCS ni en la bitácora ni en la hoja.
  \item En la hoja, tres vacas tienen la nota ``Fresca'' pero VMS 3 las marca como preñada o inseminada, que son la 1565, la 8744 y la 91576.
  \item 15 vacas tienen BCS de las dos fechas y en 4 de ellas la diferencia es mayor a 0.25.
  \item Todas las filas de la bitácora tienen carpeta de fotos del 14 de octubre y todas esas carpetas tienen registro en la bitácora. De las vacas de la hoja, la 1578 y la 91561 no tienen fotos de la primera visita, aunque sí del 14 de octubre. No hay BCS vacíos, fuera del rango de 1 a 5 ni fuera de los pasos de 0.25 en ninguna de las dos fuentes.
\end{itemize}

\subsubsection*{Calidad de los datos de ordeño}

\begin{itemize}
  \item 233 filas tienen una acción distinta de ordeño (172 rechazadas y 61 echadas sin ordeñar). Se excluyen de las estadísticas de producción, incidentes y destino.
  \item 25 registros comparten vaca y hora de inicio con otro registro, así que pueden ser duplicados. De momento no se eliminan, por lo que los conteos pueden estar inflados hasta en 25 filas.
  \item Las columnas RCS y Usuario están vacías en las 34 vacas. No interpretamos el vacío de RCS como ausencia de mastitis, sino como que no se registró.
  \item La vaca 6211 (cuarto TI) y la 6232 (cuarto TD) no tienen lecturas en un cuarto. Según el socio formador, son cuartos que no se ordeñan por un problema de salud. Se mantienen como faltantes y no se imputan.
  \item Hay vacas con pocos datos, ya que la 1497 tuvo un ordeño, la 8771 tuvo 7 y la 8775 tuvo 10. Se excluyen de las comparaciones entre vacas.
  \item Hay ordeños con 0 kg en las vacas 1554, 2078, 6232 y 8782, y con 0.08 kg en la 8712. Falta revisar si son ordeños fallidos antes de modelar.
  \item Las columnas de patada, incompleto, pezones no encontrados y AMD de pezón están vacías entre el 82\% y el 89\% de las filas. No se tratan como faltantes, sino como ausencia del evento.
  \item En la conductividad, el rango común es de 3.0 a 5.0 mS/cm y usamos 6.0 mS/cm como valor alto.
\end{itemize}

\subsubsection*{Riesgos de las decisiones de etiqueta}

Evaluamos los riesgos con la misma escala del plan de riesgos de la fase 1. La probabilidad va de 0 a 100\%, el impacto de 1 a 5 y la magnitud es la probabilidad por el impacto. Con una magnitud de 1.5 o más mitigamos el riesgo, y con una menor lo aceptamos y lo monitoreamos.

\begin{table}[H]
  \centering
  \small
  \caption{Riesgos de las decisiones de etiqueta.}
  \label{tab:riesgos}
  \begin{tabularx}{\textwidth}{>{\raggedright\arraybackslash}p{0.33\textwidth}rrr>{\raggedright\arraybackslash}X}
    \toprule
    Riesgo & P & I & M & Decisión \\
    \midrule
    BCS distinto para la misma vaca (28 vacas con dos o más evaluaciones) & 90\% & 3 & 2.7 & Mitigar conservando las evaluaciones, entrenando con el promedio y evaluando contra cada una. \\
    \addlinespace
    Usar ``Evento BCS'' como etiqueta sin saber qué mide & 60\% & 5 & 3.0 & Mitigar usando solo el BCS de los expertos hasta que el socio confirme qué mide la columna. \\
    \addlinespace
    Nota ``Fresca'' contradictoria entre la hoja y el conjunto de datos & 75\% & 3 & 2.25 & Mitigar auditando con el CAETEC las 3 vacas en conflicto (1565, 8744 y 91576). \\
    \addlinespace
    Pocas vacas por clase de BCS, lo que vuelve inestables el MAE y los porcentajes dentro de $\pm0.25$ y $\pm0.50$ & 90\% & 4 & 3.6 & Mitigar separando entrenamiento y prueba por vaca y repartiendo las clases de BCS en los pliegues. \\
    \addlinespace
    Fotos duplicadas o asignadas a la vaca equivocada & 90\% & 4 & 3.6 & Mitigar detectando duplicados por contenido y excluyendo las fotos y vacas dudosas. \\
    \bottomrule
  \end{tabularx}
\end{table}
